\documentclass[11pt]{article}
\usepackage{mathblatt}
% Präambel des Setzers (werkzeuge/setzer-praeambel.tex, wortgleich); dazu \gteil (Muster T2U)
\makeatletter
\setlength{\mbpfbe}{0mm}% keine Punkte (pruefung.md § 5: nur beim Prüfstein)
\newlength{\gEin}\setlength{\gEin}{0pt}
\let\gSchrift\relax
\renewcommand{\pfkreuzzeile}[1]{\par\noindent\hangindent3.2em\hangafter1\hspace*{1.6em}$\square$\ #1\par}
\newcommand{\gkreuz}[1]{$\square$\ #1\hspace{2em}}
\newcommand{\gkreuzzeile}[2]{\par\noindent\hangindent2.2em\hangafter1\mbox{#1\,$\square$}\ #2\par}% Buchstabe, Option: eine je Zeile (neben Skizze)
% links, was man ansieht, rechts, was man tut; Spaltengrenze überall an derselben Stelle (gemeinsam.md „Satz“)
% v0.17: Fußeintrag „Nr: Wert“ – kurz in einer Box (nie getrennt), lang umbrechbar mit der Nummer am Anfang
% (vorher \mbox: lange Lösungen liefen über den Rand, AS40 138 pt); robust, weil \fusshilfe in die .aux schreibt
\NewDocumentCommand{\gfuss}{m m}{\sbox0{#1:~#2}\ifdim\wd0<0.45\textwidth\mbox{#1:~#2}\else\mbox{#1:}\nobreak\ #2\fi}
\newsavebox{\gbox}\newlength{\gLinks}
% v0.10 (Lehrer 13.10.): keine starre Halbierung – die linke Spalte ist so breit wie die Figur (+ 6 mm), die
% Antwort steht direkt daneben; ist die Figur breiter als 0,62 Zeile, steht die Antwort darunter
\newcommand{\gzwei}[2]{\par\smallskip\noindent\sbox{\gbox}{#1}%
  \setlength{\gLinks}{\dimexpr\wd\gbox+6mm\relax}%
  \ifdim\gLinks>0.62\linewidth
    \usebox{\gbox}\par\smallskip\noindent\begin{minipage}{\linewidth}#2\end{minipage}%
  \else
    \begin{minipage}[t]{\gLinks}\vspace{0pt}\usebox{\gbox}\end{minipage}%
    \begin{minipage}[t]{\dimexpr\linewidth-\gLinks\relax}\vspace{0pt}\raggedright #2\end{minipage}%
  \fi\par}
% 16.10. (Lehrer, PA10 S. 1): Text links, Koordinatensystem rechts daneben; ist das System breiter als 0,58 Zeile,
% steht es wie bisher unter dem Text
% Passt das System knapp nicht mehr auf die Seite (bis 25 % zu hoch), wird es auf den Rest verkleinert statt
% die Seite halb leer zu lassen
\newlength{\gRest}\newlength{\gHoch}
\newcommand{\gtextks}[2]{\sbox{\gbox}{#1}%
  \setlength{\gHoch}{\dimexpr\ht\gbox+\dp\gbox\relax}%
  \ifdim\pagegoal<\maxdimen\setlength{\gRest}{\dimexpr\pagegoal-\pagetotal-4\baselineskip\relax}%
    \ifdim\gHoch>\gRest\ifdim\gRest>0.75\gHoch
      \sbox{\gbox}{\resizebox*{!}{\gRest}{\usebox{\gbox}}}\fi\fi\fi
  \ifdim\wd\gbox>0.58\linewidth #2\par\smallskip\noindent\gpassend{\usebox{\gbox}}\par
  \else\begin{minipage}[t]{\dimexpr\linewidth-\wd\gbox-5mm\relax}\raggedright #2\end{minipage}\hfill
    \begin{minipage}[t]{\wd\gbox}\vspace{0pt}\usebox{\gbox}\end{minipage}\par\fi}
% v0.9 (12.10., pruefung.md § 5 „Marke mit B, Bedingungen rechts“): drittes Argument von pfaufg = Bedingungen
% der Prüfung (v0.10: \pfohnetr = TR durchgestrichen, „FOR“), klein grau im rechten Rand auf der Grundlinie der Nummer, Zeilen mit \\
\newcommand{\gbedrand}[2]{\ifblank{#2}{}{\rlap{\hspace*{\dimexpr#1+2mm\relax}%
  \parbox[t]{15mm}{\raggedright\scriptsize\leavevmode\color{mbgrau}\strut#2}}}}
% Teil eines Bündels (a, b …): Marke links im Rand, Bedingungen rechts im Rand, auf der Zeile des Teils
\newcommand{\gteilrand}[2]{\llap{\mbpfrandmarke{#1}\hspace*{\mbpfnr}}\gbedrand{\linewidth}{#2}}
% Teil eines Bündels: Buchstabe in eigener Spalte, Inhalt eingerückt (\gEin = Einzug, damit die Spaltengrenze von
% \gzwei an derselben Stelle bleibt)
\newenvironment{gbteil}[3]{\par\smallskip\noindent\gteilrand{#2}{#3}\makebox[\mbpfteil][l]{\textbf{#1}}%
  \begin{minipage}[t]{\dimexpr\linewidth-\mbpfteil\relax}\setlength{\gEin}{\mbpfteil}\raggedright\strut\ignorespaces}%
 {\end{minipage}\par\addvspace{6pt}}% v0.18 (PG80 Nr. 1): Abstand, damit die Linie von a) nicht an b) rutscht
% Grafik höchstens so breit wie die Spalte (Bündel nebeneinander, Koordinatensystem ohne Karo)
\newsavebox{\gpbox}
\newcommand{\gpassend}[1]{\sbox{\gpbox}{#1}\ifdim\wd\gpbox>\linewidth\typeout{GPASSEND: verkleinert \the\wd\gpbox\space auf \the\linewidth}\resizebox{\linewidth}{!}{\usebox{\gpbox}}\else\usebox{\gpbox}\fi}
% Skizzenfeld (v0.10, Lehrer 13.10.): Karo ab dem linken Rand über die Zeile, „Skizze“ klein grau oben links im Feld
% (wie \pfskizzenfeld der Vorlage); der rechte Rand bleibt den Bedingungen
\newcommand{\gskizze}[1]{\par\smallskip\noindent\begin{tikzpicture}
  \clip (0,0) rectangle ({\linewidth-1mm},{#1*10mm});
  \draw[black!16,line width=0.3pt,step=5mm] (0,0) grid ({\linewidth},{#1*10mm});
  \node[anchor=north west,inner sep=2pt,font=\scriptsize,text=mbgrau,fill=white] at (0,{#1*10mm}) {Skizze};
  \end{tikzpicture}\par}
\RenewDocumentEnvironment{pfaufg}{m m m}{%
  \begin{lrbox}{\mbaufgabebox}\begin{minipage}[t]{\linewidth}%
  \gSchrift\noindent\hspace*{\gEin}\mbpfrandmarke{#1}\gbedrand{\linewidth-\gEin}{#3}%
  \begin{minipage}[t]{\mbpfnr}\strut\textbf{#2}\end{minipage}%
  \begin{minipage}[t]{\dimexpr\linewidth-\gEin-\mbpfnr\relax}\raggedright\strut\ignorespaces}%
 {\end{minipage}%
  \end{minipage}\end{lrbox}%
  \setlength{\mbaufgabehoehe}{\dimexpr\ht\mbaufgabebox+\dp\mbaufgabebox\relax}%
  \ifdefempty{\gkopfnach}{\par\addvspace{5pt}\Needspace*{\mbaufgabehoehe}}{%
    % Stufenkopf, der mit dieser (ersten) Aufgabe zusammenbleibt: Platz für beide verlangen, dann den Kopf setzen
    % (12.10., Quadratische: ein Kopf stand allein unten, die Aufgabe mit Koordinatensystem auf der nächsten Seite)
    % v0.17 (AS40 „Terme deuten“ allein unten): hinter dem Kopf steht das Sprungziel (\special, ein Knoten, nach
    % dem die Seite brechen durfte) – darum \nobreak nach dem Kopf; Vorhalt 56 pt = Kopf mit Linie und Abständen
    \par\Needspace*{\dimexpr\mbaufgabehoehe+56pt\relax}\gkopfnach\par\nobreak\gdef\gkopfnach{}\addvspace{5pt}}%
  \noindent\usebox{\mbaufgabebox}\par\addvspace{5pt}}
% v0.18 (plan.md § 10, 15.10. l): Abschnittsbuchstabe A, B, C vor dem Namen im Kopf – als Vorsatz, damit der
% Name im TeX gleich bleibt (Bestellkarte, Sprungziele lesen ihn)
\def\gbstakt{}\newcommand{\gbstvor}[1]{\gdef\gbstakt{#1}}
\NewCommandCopy\gpfstufekopfalt\pfstufekopf
\RenewDocumentCommand{\pfstufekopf}{o m m}{\IfValueTF{#1}{\gpfstufekopfalt[#1]}{\gpfstufekopfalt}%
  {\ifx\gbstakt\empty\else\gbstakt\enspace\fi #2}{#3}\gdef\gbstakt{}}
% v0.18 (16.10. e): Begriff grau, ohne Nummer, ohne Rahmen über der ersten Aufgabe der Stufe
\newcommand{\gbegriff}[1]{\par\nobreak\noindent\hspace*{\mbpfnr}\parbox[t]{\dimexpr\linewidth-\mbpfnr\relax}{\raggedright\small\color{mbgrau}#1}\par\nobreak\addvspace{3pt}}
% v0.18 (16.10. g): Lösungsblatt – Kopf mit kräftigem Strich, Abschnitt klein grau als Zwischenzeile
\newcommand{\glkopf}[1]{\noindent{\large\bfseries Lösungen \textperiodcentered{} #1}\par\vspace{3pt}\hrule height 1.2pt\par\medskip}
\newcommand{\glzwischen}[1]{\par\addvspace{8pt}\Needspace*{4\baselineskip}\noindent{\footnotesize\color{mbgrau}#1}\par\nobreak\vspace{1pt}}
% v0.18 (15.10. o, w): Blattende – Kern ohne Rahmen, dann graue Zeile „Mehr: n Aufgaben“
\newcommand{\gkern}[1]{\par\addvspace{14pt}\Needspace*{6\baselineskip}\noindent{\bfseries Kern}\par\nobreak\smallskip
  {\small\raggedright #1\par}}
\newcommand{\gmehr}[1]{\par\addvspace{10pt}\noindent{\footnotesize\color{mbgrau}Mehr: #1 Aufgaben}\par}
% Stufenkopf vormerken: \gkopfvor{<Kopf>} – die nächste pfaufg setzt ihn, sobald beide zusammen Platz haben
\def\gkopfnach{}
\newcommand{\gkopfvor}[1]{\gdef\gkopfnach{#1}}
% Teilaufgabe mit Skizze links in gzwei: Buchstabe in eigener Spalte, darunter Text und Skizze
\newcommand{\gteil}[2]{\makebox[\mbpfteil][l]{\textbf{#1}}\parbox[t]{\dimexpr0.5\textwidth-\gEin-\mbpfnr-5mm-\mbpfteil\relax}{\raggedright #2}}
% Stufenkopf der Vorlage (mathblatt.sty 2026-10-12: Linie, Abstand, 13 pt, \pfjahre) ohne ihren Vorhalt von
% 14 Zeilen: der Setzer fordert vor dem Kopf 0,16 Seitenhöhe (Kopf + erste Aufgabe), mehr nicht
\setlength{\mbpfstufevorhalt}{0pt}
\makeatother
\newcommand{\karo}[1]{\par\smallskip\noindent\begin{tikzpicture}
  \clip (0,0) rectangle ({\linewidth-1mm},{#1*10mm});
  \draw[black!16,line width=0.3pt,step=5mm] (0,0) grid ({\linewidth},{#1*10mm});
  \end{tikzpicture}\par}
\newcommand{\kk}[1]{$\square$\,#1\hspace{0.9em}}
\newcommand{\linie}[1]{\,{\color{black!45}\rule[-1pt]{#1}{0.4pt}}\,}
\newcommand{\pkt}{\textperiodcentered{}}
% v0.14 (14.10.): Stufenkopf mit fünf Punkten (Option = Jahre für die Bestellkarte) oder „zuletzt 19“
\NewDocumentCommand{\gpunkte}{o m}{#2}
\newcommand{\gzuletzt}[1]{zuletzt #1}
% v0.7 (12.10.): Schülerfrage unter dem Stufenkopf; Ankreuzoption als Bild (Buchstabe und Kästchen darunter)
\newcommand{\gfrage}[1]{\par\nobreak\noindent\hspace*{\mbpfnr}{\small\itshape #1}\par\nobreak\addvspace{2pt}}
\newlength{\gbildbreite}\setlength{\gbildbreite}{0.32\linewidth}% Setzer setzt sie je Reihe (vier Bilder: 0,24)
\newcommand{\gbildoption}[3]{\begin{minipage}[t][#3][b]{\gbildbreite}\centering #2\par\smallskip\mbox{#1\,$\square$}\end{minipage}}
% Lineare (12.10.): Bild mit Buchstabe ohne Kästchen (Diagramm zuordnen); Aussage mit „wahr / falsch“ zum Ankreuzen
\newcommand{\gbildname}[3]{\begin{minipage}[t][#3][b]{\gbildbreite}\centering #2\par\smallskip\mbox{\textbf{#1}}\end{minipage}}
\newcommand{\gwahrfalsch}[2]{\par\noindent\hangindent2.2em\hangafter1\makebox[2.2em][l]{(#1)}#2\hfill\mbox{wahr $\square$\quad falsch $\square$}\par}
\begin{document}

\pfheftstil
\pflaufkopf[P10 \pkt{} Prüfung]{Brüche und Potenzen}{PZ10-2}\pfseitevon
\pfheftkopf{Brüche und Potenzen}{}
\pfstufenzeile{A~Anteil gesucht \pkt{} B~Zahlen vergleichen \pkt{} C~Potenzen umwandeln}

% ===================== Stufe: Anteil gesucht =====================
\gkopfvor{\gbstvor{A}\pfstufekopf{Anteil gesucht}{\gpunkte[26 25 24 23 22]{\textcolor{black!80}{$\bullet$}\,\textcolor{black!80}{$\bullet$}\,\textcolor{black!80}{$\bullet$}\,\textcolor{black!80}{$\bullet$}\,\textcolor{black!80}{$\bullet$}}}\special{pdf:dest (stufe-anteil-gesucht) [@thispage /XYZ @xpos @ypos null]}}
% kein Einstieg: keine Schülerfrage in der Gliederung (Schülerfrage, -Fälle, -Form, -Weg, -Skizze; oder Frage) – Einstieg nur von dort, nie aus der Bank (pruefung.md § 5, 12.10.)
% 2022-OS-B1a-m1 (Original 2022-OS-B1a, Schwere I, rang 1; herausgelöste Fassung: keine)
\begin{pfaufg}{’22~\pkt{}~1a}{1.}{Basis}
Markiere 25 \% der Fläche.
\par\smallskip\noindent\begin{tikzpicture}[x=0.45cm,y=0.45cm,line width=0.5pt]\draw[line width=0.4pt] (1,0) -- (1,4);\draw[line width=0.4pt] (2,0) -- (2,4);\draw[line width=0.4pt] (3,0) -- (3,4);\draw[line width=0.4pt] (4,0) -- (4,4);\draw[line width=0.4pt] (0,1) -- (5,1);\draw[line width=0.4pt] (0,2) -- (5,2);\draw[line width=0.4pt] (0,3) -- (5,3);\draw (0,0) rectangle (5,4);\end{tikzpicture}\par
\fusshilfe{\gfuss{1}{5 Kästchen}}
\end{pfaufg}
% 2023-OS-B1d-m1 (Original 2023-OS-B1d, Schwere I, rang 2; herausgelöste Fassung: keine)
\begin{pfaufg}{’23~\pkt{}~1d}{2.}{Basis}
Markiere 30 \% der Fläche.
\par\smallskip\noindent\begin{tikzpicture}[x=0.35cm,y=0.35cm,line width=0.5pt]\draw (2,0) rectangle ++(1,1);\draw (3,0) rectangle ++(1,1);\draw (1,1) rectangle ++(1,1);\draw (2,1) rectangle ++(1,1);\draw (3,1) rectangle ++(1,1);\draw (4,1) rectangle ++(1,1);\draw (0,2) rectangle ++(1,1);\draw (1,2) rectangle ++(1,1);\draw (2,2) rectangle ++(1,1);\draw (3,2) rectangle ++(1,1);\draw (4,2) rectangle ++(1,1);\draw (5,2) rectangle ++(1,1);\draw (0,3) rectangle ++(1,1);\draw (1,3) rectangle ++(1,1);\draw (2,3) rectangle ++(1,1);\draw (3,3) rectangle ++(1,1);\draw (4,3) rectangle ++(1,1);\draw (5,3) rectangle ++(1,1);\draw (0,4) rectangle ++(1,1);\draw (1,4) rectangle ++(1,1);\draw (2,4) rectangle ++(1,1);\draw (3,4) rectangle ++(1,1);\draw (4,4) rectangle ++(1,1);\draw (5,4) rectangle ++(1,1);\draw (1,5) rectangle ++(1,1);\draw (2,5) rectangle ++(1,1);\draw (3,5) rectangle ++(1,1);\draw (4,5) rectangle ++(1,1);\draw (2,6) rectangle ++(1,1);\draw (3,6) rectangle ++(1,1);\end{tikzpicture}\par
\fusshilfe{\gfuss{2}{9 Kästchen}}
\end{pfaufg}
% 2024-OS-B1b-m1 (Original 2024-OS-B1b, Schwere I, rang 3; herausgelöste Fassung: keine)
\begin{pfaufg}{’24~\pkt{}~1b}{3.}{Basis}
Von dem Quadrat wurde ein Teil grau gefärbt. Gib den Anteil der grauen Fläche an.
\gzwei{\begin{tikzpicture}[line width=0.6pt]\fill[black!25] (0,0) -- (2,0) -- (1,1) -- cycle;\draw[line width=0.3pt,step=0.5] (0,0) grid (2,2);\draw (0,0) -- (1,1) -- (2,0);\draw (0,0) rectangle (2,2);\end{tikzpicture}}{\vspace{1mm}\par\smallskip\par\noindent \hfill Anteil: \linie{30mm}\hphantom{\,cm²}}
\fusshilfe{\gfuss{3}{25 \%}}
\end{pfaufg}
% 2025-OS-B1c-m1 (Original 2025-OS-B1c, Schwere I, rang 4; herausgelöste Fassung: keine)
\begin{pfaufg}{’25~\pkt{}~1c}{4.}{Basis}
Kennzeichne in dem Rechteck $\tfrac{1}{3}$ der Fläche.
\par\smallskip\noindent\begin{tikzpicture}[line width=0.6pt]\draw (0,0) rectangle (3,1.5);\end{tikzpicture}\par
\fusshilfe{\gfuss{4}{Streifen 1 cm von 3 cm breit}}
\end{pfaufg}
% 2026-FOR-B1b-m1 (Original 2026-FOR-B1b, Schwere I, rang 5; herausgelöste Fassung: keine)
\begin{pfaufg}{’26~\pkt{}~1b}{5.}{Basis}
Kreuze an, in welcher Abbildung die graue Fläche einem Anteil von $\tfrac{1}{4}$ entspricht.
\par\smallskip\noindent\setlength{\gbildbreite}{\dimexpr0.96\linewidth/4\relax}\sbox{\gbox}{\vbox{\begin{tikzpicture}[line width=0.6pt]\fill[black!40] (0,0) rectangle (1,1.5);\draw (0,0) rectangle (3,1.5);\draw (1,0) -- (1,1.5);\end{tikzpicture}\par\begin{tikzpicture}[line width=0.6pt,rotate=45]\fill[black!40] (0,0) rectangle (0.75,0.75);\draw (0,0) rectangle (1.5,1.5);\draw (0.75,0) -- (0.75,1.5) (0,0.75) -- (1.5,0.75);\end{tikzpicture}\par\begin{tikzpicture}[line width=0.6pt]\fill[black!40] (0,0) -- (0.55,1.25) -- (2.25,1.25) -- cycle;\draw (0,0) -- (0.55,1.25) -- (2.25,1.25) -- (1.7,0) -- cycle;\draw (0,0) -- (2.25,1.25);\end{tikzpicture}\par\begin{tikzpicture}[line width=0.6pt]\filldraw[fill=black!40] (0,0) -- (90:0.78) arc (90:150:0.78) -- cycle;\draw (0,0) circle (0.78);\end{tikzpicture}}}\gbildoption{}{\begin{tikzpicture}[line width=0.6pt]\fill[black!40] (0,0) rectangle (1,1.5);\draw (0,0) rectangle (3,1.5);\draw (1,0) -- (1,1.5);\end{tikzpicture}}{\dimexpr\ht\gbox/4+\baselineskip\relax}\hfill\gbildoption{}{\begin{tikzpicture}[line width=0.6pt,rotate=45]\fill[black!40] (0,0) rectangle (0.75,0.75);\draw (0,0) rectangle (1.5,1.5);\draw (0.75,0) -- (0.75,1.5) (0,0.75) -- (1.5,0.75);\end{tikzpicture}}{\dimexpr\ht\gbox/4+\baselineskip\relax}\hfill\gbildoption{}{\begin{tikzpicture}[line width=0.6pt]\fill[black!40] (0,0) -- (0.55,1.25) -- (2.25,1.25) -- cycle;\draw (0,0) -- (0.55,1.25) -- (2.25,1.25) -- (1.7,0) -- cycle;\draw (0,0) -- (2.25,1.25);\end{tikzpicture}}{\dimexpr\ht\gbox/4+\baselineskip\relax}\hfill\gbildoption{}{\begin{tikzpicture}[line width=0.6pt]\filldraw[fill=black!40] (0,0) -- (90:0.78) arc (90:150:0.78) -- cycle;\draw (0,0) circle (0.78);\end{tikzpicture}}{\dimexpr\ht\gbox/4+\baselineskip\relax}\par
\end{pfaufg}

% ===================== Stufe: Zahlen vergleichen =====================
\gkopfvor{\gbstvor{B}\pfstufekopf{Zahlen vergleichen}{\gpunkte[23 22]{\textcolor{black!25}{$\circ$}\,\textcolor{black!25}{$\circ$}\,\textcolor{black!25}{$\circ$}\,\textcolor{black!80}{$\bullet$}\,\textcolor{black!80}{$\bullet$}}}\special{pdf:dest (stufe-zahlen-vergleichen) [@thispage /XYZ @xpos @ypos null]}}
% kein Einstieg: keine Schülerfrage in der Gliederung (Schülerfrage, -Fälle, -Form, -Weg, -Skizze; oder Frage) – Einstieg nur von dort, nie aus der Bank (pruefung.md § 5, 12.10.)
% 2023-OS-B1f-m1 (Original 2023-OS-B1f, Schwere I, rang 12; herausgelöste Fassung: keine)
\begin{pfaufg}{’23~\pkt{}~1f}{6.}{Basis}
Gib den kleinsten Wert an: 3,3 · 0,33 · $0{,}3^2$ · 33 \%
\hfill\linie{30mm}
\fusshilfe{\gfuss{6}{0,09}}
\end{pfaufg}
% 2022-OS-B1j-m1 (Original 2022-OS-B1j, Schwere I, rang 13; herausgelöste Fassung: keine)
\begin{pfaufg}{’22~\pkt{}~1j}{7.}{Basis}
Setze das richtige Zeichen (<, = oder >) ein: $2^{ - 3}$ \linie{25mm} 0,1
\fusshilfe{\gfuss{7}{0,125 > 0,1, also >}}
\end{pfaufg}

% ===================== Stufe: Potenzen umwandeln =====================
\gkopfvor{\gbstvor{C}\pfstufekopf{Potenzen umwandeln}{\gpunkte[25 24 22]{\textcolor{black!25}{$\circ$}\,\textcolor{black!80}{$\bullet$}\,\textcolor{black!80}{$\bullet$}\,\textcolor{black!25}{$\circ$}\,\textcolor{black!80}{$\bullet$}}}\special{pdf:dest (stufe-potenzen-umwandeln) [@thispage /XYZ @xpos @ypos null]}}
% kein Einstieg: keine Schülerfrage in der Gliederung (Schülerfrage, -Fälle, -Form, -Weg, -Skizze; oder Frage) – Einstieg nur von dort, nie aus der Bank (pruefung.md § 5, 12.10.)
% 2025-OS-B1d-m1 (Original 2025-OS-B1d, Schwere I, rang 21; herausgelöste Fassung: keine)
\begin{pfaufg}{’25~\pkt{}~1d}{8.}{Basis}
Gib den Exponenten so an, dass die Gleichung stimmt: 3 600 000 = 3,6 · $10^{\square}$
\par\smallskip\par\noindent \hfill Exponent: \linie{30mm}\hphantom{\,cm²}
\fusshilfe{\gfuss{8}{Exponent 6}}
\end{pfaufg}
% 2024-OS-B1g-m1 (Original 2024-OS-B1g, Schwere I, rang 22; herausgelöste Fassung: keine)
\begin{pfaufg}{’24~\pkt{}~1g}{9.}{Basis}
Gib den Wert für x an, sodass gilt: $5^x = 125$
\hfill\linie{30mm}
\fusshilfe{\gfuss{9}{3}}
\end{pfaufg}

\gmehr{18}
\end{document}